\chapter[Le fichier séquentiel]{
Le fichier séquentiel}
{
Dans ce chapitre, on voit comment manipuler les fichiers séquentiels via
deux notations : une notation procédurale et une notation orienté objet
(nous vous laisserons utiliser celle que vous préférez). }

\begin{center}
 [Warning: Image ignored] % Unhandled or unsupported graphics:
%\includegraphics[width=1.27cm,height=1.27cm]{log1-img/log1-img157}

\end{center}
{
Ensuite, on généralise la notion de fichier à la notion de «~séquence~»
qui peut englober aussi d'autres structures comme les
listes et les tableaux.}

{
Dans la première partie de ce chapitre, nous utilisons la convention
courante qu'une lecture infructueuse est nécessaire
pour détecter la fin d'un fichier. En fin de chapitre
on voit une autre convention de lecture qu'on
rencontre dans certains langages et qui stipule qu'on
peut détecter la fin de fichier juste après la dernière lecture utile.
}

\section{Présentation}
{
Un \textbf{fichier} informatique est une unité informationnelle
physiquement stockée sur un support de mémoire de masse permanent
(disque dur, CD-ROM, clé USB par exemple). }

{
Le qualificatif \textbf{séquentiel}\footnote{Il existe également des
fichiers à \textit{accès direct} (cf. votre cours
d'organisation des données)} signifie que
lorsqu'on \textit{lit} un fichier, les éléments sont
accédés dans une \textbf{séquence préétablie et ordonnée} (pour accéder
à un élément de position donnée, il faut d’abord avoir \textit{lu} tous
les éléments précédents). De même, lorsqu'on
\textit{écrit} un fichier séquentiel, les éléments sont ordonnés dans
l'ordre dans lequel ils sont écrits (une prochaine
lecture les donnerait dans ce même ordre).}

{
Nous allons considérer qu'un fichier peut être lu ou
écrit mais pas les deux en même temps\footnote{Certains systèmes et
langages mettent à disposition des fichiers en lecture/écriture
simultanées}. On considère aussi qu'un fichier doit
explicitement être \textbf{ouvert} avant toute lecture/écriture et
qu'il doit être \textbf{fermé} après la dernière
lecture/écriture.}

{
Lors d'une lecture de fichier, on appelle
«~\textbf{élément courant~»}, le prochain élément qui sera lu.
L’ensemble étant fini, il doit être possible d’en détecter la
\textbf{fin}; on peut imaginer la présence d’une marque particulière
caractérisant cette fin (\textbf{marque de fin de fichier}). }

{
On comprend donc qu'il faut une lecture supplémentaire
explicite pour lire la fin de fichier et se rendre ainsi compte
qu'il n'y a plus rien à lire. En fin
de chapitre on présente un autre type d'accès où on
sait qu'on est à la fin du fichier juste après la
lecture du dernier élément.}

{
Le fichier n'est pas la seule structure qui peut être
accédée en séquence. On pourrait faire de même pour un tableau, une
liste ou encore d'autres structures que vous verrez en
deuxième année. Ce point est, lui aussi, abordé en fin de chapitre.}

\section{Primitives – notation procédurale}
{
Voyons les notations utilisées pour manipuler des fichiers dans un
contexte non orienté-objet.}

{\sffamily\bfseries\upshape
Déclaration d’un fichier séquentiel }

{
Les fichiers utilisés dans un algorithme seront déclarés dans la partie
déclarative du code selon le modèle : }

{\sffamily
nomFichier : fichier de T }

{
où \textstyleCodeInsr{T} est le type des éléments du fichier (types
simples, variables structurées, objets)}

{
Exemples:}

{\sffamily
fileNumbers : fichier d’entiers }

{\sffamily
fichierTexte : fichier de chaines }

{\sffamily
fichCalendrier : fichier de Date }

{\sffamily\bfseries\upshape
Ouverture d’un fichier en lecture }

{\sffamily
\textstyleMotCl{ouvrir} nomFichier (\textstyleMotCl{lecture}) \ // ou
\textbf{IN }ou\textbf{ INPUT} }

{
Cette primitive met à la disposition de l’exécutant le fichier
\textstyleCodeInsr{nomFichier} en spécifiant que ce fichier sera
utilisé en lecture (IN). De plus, elle positionne la tête de lecture
devant le premier élément du fichier (qui n’est donc pas encore lu).
Attention, si le fichier \textstyleCodeInsr{nomFichier} est inexistant,
cette primitive provoque l’interruption du programme. }

{
Pour ouvrir plusieurs fichiers en lecture :}

{\sffamily
\textstyleMotCl{ouvrir} nomFichier1, nomFichier2, nomFichier3
(\textstyleMotCl{lecture})\ \ // ou \textbf{IN }ou\textbf{ INPUT}}

{\sffamily\bfseries\upshape
Ouverture d’un fichier en écriture }

{\sffamily
\textstyleMotCl{ouvrir} nomFichier (\textstyleMotCl{écriture}) \ \ // ou
\textstyleMotCl{OUT} \ ou \textstyleMotCl{OUTPUT}}

{
Cette primitive met à la disposition de l’exécutant le fichier
\textstyleCodeInsr{nomFichier} en spécifiant que ce fichier sera
utilisé en écriture (OUT). De plus, elle positionne la tête d’écriture
en début du fichier. Si le fichier \textstyleCodeInsr{nomFichier} est
inexistant, celui-ci sera créé lors de l’exécution de cette primitive.
Par contre, si ce fichier existe déjà, son contenu est «~détruit~» lors
de l'ouverture.}

{
Pour ouvrir plusieurs fichiers en écriture :}

{\sffamily
\textstyleMotCl{ouvrir} nomFichier1, nomFichier2, nomFichier3
(\textstyleMotCl{écriture}) \ \ }

{\sffamily
\ \ \ \ \ \ \ \ \ \ \ \ \ \ \ \ \ \ // ou \textstyleMotCl{OUT} \ ou
\textstyleMotCl{OUTPUT}}

{\sffamily\bfseries\upshape
Lecture dans un fichier (ouvert en lecture) }

{\sffamily
\textstyleMotCl{lire} nomFichier ; nomVariable }

{
Cette primitive lit dans le fichier \textstyleCodeInsr{nomFichier} une
valeur qui sera affectée à la variable \textstyleCodeInsr{nomVariable}.
La valeur lue est la donnée du fichier se trouvant en face de la tête
de lecture. }

{
Si \textstyleCodeInsr{nomFichier} est un fichier de T alors
\textstyleCodeInsr{nomVariable} doit être déclaré de type T.}

{
Pour détecter la fin de fichier, on utilise la
\textbf{notation}\footnote{Dans certains langages, il
s'agira d'une variable en lecture
seule; dans d'autres langages, ce sera un appel de
module ou de méthode.}\textbf{ }\textstyleCodeInsr{EOF(nomFichier)}.
Après une lecture \textstyleCodeInsr{EOF(nomFichier)} prend la valeur
\textbf{vrai} si la \textbf{fin de fichier} a été détectée et
\textbf{faux }sinon. }

{
Dans le cas où \textstyleCodeInsr{EOF(nomFichier)} est à
\textbf{vrai}\textstyleCodeInsr{,} le contenu de
\textstyleCodeInsr{nomVariable} n’est plus affecté et doit être
considéré comme indéterminé. }

{
Attention à ne pas commettre les erreurs d’utilisation suivantes qui
provoquent toutes l’interruption de l’algorithme : }

\liststyleListv
\begin{itemize}
\item {
lecture dans un fichier ouvert en écriture, ou non ouvert en lecture; }
\item {
utilisation de la primitive \textstyleMotCl{lire} alors que la marque de
fin de fichier a déjà été détectée, c’est-à-dire lorsque
\textstyleCodeInsr{EOF(nomFichier)} est devenu \textbf{vrai};}
\item {
utilisation de \textstyleCodeInsr{nomVariable} alors que
\textstyleCodeInsr{EOF(nomFichier)} est devenu \textbf{vrai};}
\item {
appel à \textstyleCodeInsr{EOF(nomFichier)} alors
qu'aucune lecture n'a encore été
faite.}
\end{itemize}
{\sffamily\bfseries
Écriture dans un fichier (ouvert en écriture) }

{\sffamily
\textstyleMotCl{écrire} nomFichier ; nomVariable }

{
Cette primitive réalise l’écriture du contenu de la variable
\textstyleCodeInsr{nomVariable} dans le fichier
\textstyleCodeInsr{nomFichier} à l’endroit correspondant à la position
de la tête d’écriture. Après l’écriture, la tête avance d’une position.
Il est bien entendu interdit d’écrire dans un fichier ouvert en
lecture, ou qui n’a pas été ouvert en écriture. }

{\sffamily\bfseries\upshape
Fermeture d’un fichier }

{\sffamily
\textstyleMotCl{fermer} nomFichier }

{
Cette primitive permet de fermer un fichier après son utilisation, en
lecture ou en écriture. Dans le cas où le fichier a été ouvert en
écriture, cette primitive place la marque spéciale de fin de fichier
dans l’élément courant. Une fois le fichier fermé, il
n'est plus permis de l'utiliser.}

{
Pour fermer plusieurs fichiers :}

{\sffamily
\textstyleMotCl{fermer} nomFichier1, nomFichier2, nomFichier3 }

\section{Algorithme de base }
{
L’algorithme suivant parcourt simplement le fichier
\textstyleCodeInsr{nomFichier}. Il s’agit de l’algorithme type dont la
structure se retrouve dans la majorité des algorithmes exploitant un
fichier. }

{\sffamily
\textstyleMotCl{module} parcoursFichier()}

{\sffamily
\ \ nomFichier : fichier de T }

{\sffamily
\ \ enregistrement : T }

{\sffamily
\ \ \textstyleMotCl{ouvrir} nomFichier (lecture) }

{\sffamily
\ \ \textstyleMotCl{lire} nomFichier ; enregistrement }

{\sffamily
\ \ \textstyleMotCl{tant} \textstyleMotCl{que} NON EOF(nomFichier)
\textstyleMotCl{faire} }

{\sffamily\itshape
\ \ \ \ // traitement de l’information contenue dans enregistrement}

{\sffamily
\ \ \ \ \textstyleMotCl{lire} nomFichier ; enregistrement }

{\sffamily
\ \ \textstyleMotCl{fin} \textstyleMotCl{tant} \textstyleMotCl{que} }

{\sffamily
\ \ \textstyleMotCl{fermer} nomFichier }

{\sffamily
\textstyleMotCl{fin module}}


\bigskip

{\sffamily\bfseries\scshape
Exemple : Copie d'un fichier}

{
Par exemple, voici un programme qui recopie simplement les données
d'un fichier (en lecture) dans un autre (en écriture)
}

{\sffamily
\textstyleMotCl{module} copieFichier() }

{\sffamily
\ \ fichierIn : fichier de T }

{\sffamily
\ \ fichierOut : fichier de T }

{\sffamily
\ \ enregistrement : T }

{\sffamily
\ \ \textstyleMotCl{ouvrir} fichierIn (lecture), \ fichierOut
(écriture)}

{\sffamily
\ \ \textstyleMotCl{lire} fichierIn ; enregistrement }

{\sffamily
\ \ \textstyleMotCl{tant} \textstyleMotCl{que} NON EOF(fichierIn)
\textstyleMotCl{faire} }

{\sffamily
\ \ \ \ \textstyleMotCl{écrire} fichierOut ; enregistrement}

{\sffamily
\ \ \ \ \textstyleMotCl{lire} fichierIn ; enregistrement }

{\sffamily
\ \ \textstyleMotCl{fin} \textstyleMotCl{tant} \textstyleMotCl{que} }

{\sffamily
\ \ \textstyleMotCl{fermer} fichierIn, fichierOut }

{\sffamily
\textstyleMotCl{fin module}}

\section{Primitives – notation orienté objet}
{
On peut aussi définir une notation orienté objet pour la manipulation
d'un fichier. On utilise une classe pour les fichiers
séquentiels en entrée et une autre pour les fichiers séquentiels en
sortie.}

{\sffamily
\textstyleMotCl{classe} FichierEntrée de T}

{\sffamily
\textstyleMotCl{public} :}

{\sffamily
\ \ \textstyleMotCl{constructeur} FichierEntrée de T(localisation :
chaine)}

{\sffamily
\ \ \textstyleMotCl{méthode} ouvrir()}

{\sffamily
\ \ \textstyleMotCl{méthode} lire() {→} T}

{\sffamily
\ \ \textstyleMotCl{méthode} EOF() {→} booléen}

{\sffamily
\ \ \textstyleMotCl{méthode} fermer()}

{\sffamily
\textstyleMotCl{fin} \textstyleMotCl{classe}}


\bigskip

{\sffamily
\textstyleMotCl{classe} FichierSortie de T}

{\sffamily
\textstyleMotCl{public} :}

{\sffamily
\ \ \textstyleMotCl{constructeur} FichierSortie de T(localisation :
chaine)}

{\sffamily
\ \ \textstyleMotCl{méthode} ouvrir()}

{\sffamily
\ \ \textstyleMotCl{méthode} écrire( enr : T )}

{\sffamily
\ \ \textstyleMotCl{méthode} fermer()}

{\sffamily
\textstyleMotCl{fin }\textstyleMotCl{classe}}

{
Quelques remarques s'imposent :}

\liststyleListv
\begin{itemize}
\item {
Dans la version non orienté objet, rien n'est dit sur
la localisation du fichier; on suppose que cela est déterminé par
ailleurs. Ici, nous supposerons que le lien est fait à la
construction.}
\item {
En Java, l'appel au constructeur «~ouvre~» le fichier;
la tête de lecture/écriture est positionnée au début. Toutefois, pour
rester proche de la version non OO, nous supposerons en logique
qu'un appel à la méthode \textstyleCodeInsr{ouvrir()}
est nécessaire.}
\end{itemize}
{\sffamily\bfseries\scshape
Exemple : parcours d'un fichier}

{
Récrivons l'algorithme de base en notation orienté
objet.}

{\sffamily
\textstyleMotCl{module} parcoursFichier() }

{\sffamily
\ \ fichier : FichierEntrée de T}

{\sffamily
\ \ enregistrement : T}

{\sffamily
\ \ fichier {\textrm{[F0DF?]}} nouveau
FichierEntrée de T( {\textquotedbl}data.txt{\textquotedbl} )}

{\sffamily
\ \ fichier.ouvrir() }

{\sffamily
\ \ enregistrement {\textrm{[F0DF?]}}
fichier.lire()}

{\sffamily
\ \ \textstyleMotCl{tant} \textstyleMotCl{que} NON fichier.EOF()
\textstyleMotCl{faire} }

{\sffamily\itshape
\ \ \ \ // traitement de l’information contenue dans enregistrement}

{\sffamily
\ \ \ \ enregistrement {\textrm{[F0DF?]}}
fichier.lire()}

{\sffamily
\ \ \textstyleMotCl{fin} \textstyleMotCl{tant} \textstyleMotCl{que} }

{\sffamily
\ \ fichier.fermer() }

{\sffamily
\textstyleMotCl{fin module}}

{
Bien souvent, le fichier sera un paramètre du module. Dans ce cas, il
aura déjà été créé. Illustrons ce point pour
l'algorithme de copie de fichier.}

{\sffamily
\textstyleMotCl{module} copieFichier( fichierIn : FichierEntrée de T, }

{\sffamily
\ \ \ \ \ \ \ \  \ \ \ \ fichierOut : FichierSortie de T ) }

{\sffamily
\ \ enregistrement : T }

{\sffamily
\ \ fichierIn.ouvrir() }

{\sffamily
\ \ fichierOut.ouvrir() }

{\sffamily
\ \ enregistrement {\textrm{[F0DF?]}}
fichierIn.lire()}

{\sffamily
\ \ \textstyleMotCl{tant} \textstyleMotCl{que} NON fichierIn.EOF()
\textstyleMotCl{faire} }

{\sffamily
\ \ \ \ fichierOut.écrire( enregistrement )}

{\sffamily
\ \ \ \ enregistrement {\textrm{[F0DF?]}}
fichierIn.lire()}

{\sffamily
\ \ \textstyleMotCl{fin} \textstyleMotCl{tant} \textstyleMotCl{que} }

{\sffamily
\ \ fichierIn.fermer() }

{\sffamily
\ \ fichierOut.fermer() }

{\sffamily
\textstyleMotCl{fin module}}

\section{Le cas des fichiers vides}
{
Dans tout algorithme de manipulation de fichiers, il faut se demander ce
qu'il advient du cas d'un fichier
vide (rappelons que c'est la première lecture qui
permet d'identifier ce cas). Il y a 3 cas de figure :}

\liststyleNumberingii
\begin{enumerate}
\item {
Le fichier vide est un cas particulier tout à fait acceptable et qui est
traité correctement par l'algorithme général.}
\item {
Le fichier vide demande un traitement spécial qui n'est
pas couvert par le cas général.}
\item {
Le fichier vide dénote un problème mal posé qui ne peut être résolu.}
\end{enumerate}
{\sffamily\bfseries\scshape
Exemple 1 – Des informations à traiter}

{
Le module \textstyleCodeInsr{copierFichier} ci-dessus est un exemple du
premier cas. En cas de fichier vide, l'algorithme
fonctionne parfaitement puisqu'il créera bien un
nouveau fichier vide. De façon générale, le fichier vide peut
simplement indiquer qu'il n'y a rien
à traiter pour le moment. Par exemple, on pourrait imaginer un fichier
contenant des informations sur les étudiants qui se sont mis en ordre
d'inscription aujourd'hui et qui
serait utilisé toutes les nuits par un programme créant les comptes
informatiques pour ces étudiants. \ }

{\sffamily\bfseries\scshape
Exemple 2 – Calcul de la moyenne}

{
Si le problème de départ est de calculer la somme, la moyenne ou encore
le maximum des nombres contenus dans un fichier, alors la solution ne
peut être donnée dans le cas d'un fichier vide (cas 3)
à moins que l'analyse n'ait
clairement indiqué la réponse à donner dans ce cas (cas 2).
Concentrons-nous sur le calcul de la moyenne d’une série de nombres
contenus dans un fichier. Voici un premier exemple avec la notation OO
:}

{\sffamily
\textstyleMotCl{module} calculMoyenne( fileNb : fichierEntrée de réels )
{→} réel}


\bigskip

{\sffamily
\ \ somme, nombre : réels}

{\sffamily
\ \ cpt : entier}


\bigskip

{\sffamily
\ \ fileNb.ouvrir()}

{\sffamily
\ \ cpt {\textrm{[F0DF?]}} 0}

{\sffamily
\ \ somme {\textrm{[F0DF?]}} 0}

{\sffamily
\ \ nombre {\textrm{[F0DF?]}} fileNb.lire()}

{\sffamily
\ \ \textstyleMotCl{tant} \textstyleMotCl{que} NON fileNb.EOF()
\textstyleMotCl{faire}}

{\sffamily
\ \ \ \ cpt {\textrm{[F0DF?]}} cpt + 1}

{\sffamily
\ \ \ \ somme {\textrm{[F0DF?]}} somme + nombre}

{\sffamily
\ \ \ \ nombre {\textrm{[F0DF?]}} fileNb.lire()}

{\sffamily
\ \ \textstyleMotCl{fin} \textstyleMotCl{tant} \textstyleMotCl{que}}

{\sffamily
\ \ fileNb.fermer()}

{\sffamily
\ \ \textstyleMotCl{retourner}{\textrm{
}}{somme} / cpt}


\bigskip

{\sffamily
\textstyleMotCl{fin} \textstyleMotCl{module}}

{
Il est facile de deviner ce qui se passerait si le fichier était vide !
Le nombre de nombres du fichier étant nul, le calcul de la moyenne
provoquerait une interruption (division par zéro). }

{
Une possibilité consiste à discriminer ce cas en fin de traitement :}

{\sffamily
\textstyleMotCl{module} calculMoyenne ( fileNb : fichierEntrée de réels
) {→} réel}


\bigskip

{\sffamily
\ \ somme, nombre : réels}

{\sffamily
\ \ cpt : entier}


\bigskip

{\sffamily
\ \ fileNb.ouvrir()}

{\sffamily
\ \ cpt {\textrm{[F0DF?]}} 0}

{\sffamily
\ \ somme {\textrm{[F0DF?]}} 0}

{\sffamily
\ \ nombre {\textrm{[F0DF?]}} fileNb.lire()}

{\sffamily
\ \ \textstyleMotCl{tant} \textstyleMotCl{que} NON fileNb.EOF()
\textstyleMotCl{faire}}

{\sffamily
\ \ \ \ cpt {\textrm{[F0DF?]}} cpt + 1}

{\sffamily
\ \ \ \ somme {\textrm{[F0DF?]}} somme + nombre}

{\sffamily
\ \ \ \ nombre {\textrm{[F0DF?]}} fileNb.lire()}

{\sffamily
\ \ \textstyleMotCl{fin} \textstyleMotCl{tant} \textstyleMotCl{que}}

{\sffamily
\ \ fileNb.fermer()}

{\sffamily
\ \ \textstyleMotCl{si} cpt {\textgreater} 0 \textstyleMotCl{alors}}

{\sffamily
\ \ \ \ \textstyleMotCl{retourner}{\textrm{
}}{somme} / cpt}

{\sffamily
\ \ \textstyleMotCl{sinon}}

{\sffamily
\ \ \ \ \textstyleMotCl{erreur} {\textquotedbl}Fichier
vide{\textquotedbl}}

{\sffamily
\ \ \textstyleMotCl{fin} \textstyleMotCl{si}}

{\sffamily
\textstyleMotCl{fin} \textstyleMotCl{module}}

{
Mais ce n'est pas toujours possible. Par exemple, si le
problème est simplement de calculer la somme des nombres, la variable
\textstyleCodeInsr{cpt} disparait et la somme seule ne permet pas de
juger si le fichier est vide. En effet, la somme pourrait être nulle
parce que les différents éléments additionnés amènent précisément à ce
résultat. Il est alors plus sûr de traiter ce cas juste après la
première lecture.}

{\sffamily
\textstyleMotCl{module} calculMoyenne( fileNb : fichierEntrée de réels )
{→} réel}


\bigskip

{\sffamily
\ \ somme, nombre : réels}

{\sffamily
\ \ cpt : entier}


\bigskip

{\sffamily
\ \ fileNb.ouvrir()}

{\sffamily
\ \ cpt {\textrm{[F0DF?]}} 0}

{\sffamily
\ \ somme {\textrm{[F0DF?]}} 0}

{\sffamily
\ \ nombre {\textrm{[F0DF?]}} fileNb.lire()}


\bigskip

{\sffamily
\ \ \textstyleMotCl{si} fileNb.EOF() \textstyleMotCl{alors}}

{\sffamily
\textstyleMotCl{\ \ \ \ }\textstyleMotCl{\textmd{fileNb.fermer()}}}

{\sffamily
\ \ \ \ \textstyleMotCl{erreur} {\textquotedbl}Fichier
vide{\textquotedbl}}

{\sffamily
\ \ \textstyleMotCl{fin} \textstyleMotCl{si}}


\bigskip

{\sffamily
\ \ \textstyleMotCl{tant} \textstyleMotCl{que} \ NON fileNb.EOF()
\textstyleMotCl{faire}}

{\sffamily
\ \ \ \ cpt {\textrm{[F0DF?]}} cpt + 1}

{\sffamily
\ \ \ \ somme {\textrm{[F0DF?]}} somme + nombre}

{\sffamily
\ \ \ \ nombre {\textrm{[F0DF?]}} fileNb.lire()}

{\sffamily
\ \ \textstyleMotCl{fin} \textstyleMotCl{tant} \textstyleMotCl{que}}

{\sffamily
\ \ fileNb.fermer()}

{\sffamily
\ \ \textstyleMotCl{retourner}{\textrm{
}}{somme} / cpt}

{\sffamily
\textstyleMotCl{fin} \textstyleMotCl{module}}


\bigskip

{
Écrivons l'algorithme général de traitement
d'un fichier séquentiel :}

{\sffamily
\textstyleMotCl{module} algorithmeDeTraitementSurFichier}


\bigskip

{\sffamily
\ \ // \textit{déclaration des variables, fichiers,... éventuellement en
paramètres}}


\bigskip

{\sffamily
\ \ \textstyleMotCl{ouvrir} fichier (lecture)}

{\sffamily
\ \ \textstyleMotCl{lire} fichier ; élément}

{\sffamily
\ \ \textstyleMotCl{si} EOF(fichier) \textstyleMotCl{alors}}

{\sffamily\itshape
\ \ \ \ // traitement spécial fichier vide}

{\sffamily
\ \ \textstyleMotCl{sinon}}

{\sffamily\itshape
\ \ \ \ // initialisation du traitement de l’ensemble}

{\sffamily
\ \ \ \ \textstyleMotCl{tant} \textstyleMotCl{que} NON EOF(fichier)
\textstyleMotCl{faire}}

{\sffamily\itshape
\ \ \ \ \ \ // traitement d’un élément}

{\sffamily
\ \ \ \ \ \ \textstyleMotCl{lire} fichier ; élément}

{\sffamily
\ \ \ \ \textstyleMotCl{fin} \textstyleMotCl{tant} \textstyleMotCl{que}}

{\sffamily\itshape
\ \ \ \ // clôture du traitement de l’ensemble}

{\sffamily
\ \ \textstyleMotCl{fin} \textstyleMotCl{si}}

{\sffamily
\ \ \textstyleMotCl{fermer} fichier}

{\sffamily
\textstyleMotCl{fin} \textstyleMotCl{module}}

{
Lorsque le «~traitement spécial du fichier vide~» consiste à générer une
«~ERREUR~» alors le \textstyleMotCl{si-sinon-fin si} peut être changé
en un simple \textstyleMotCl{si-fin si} suivi du traitement normal.}

\section{La séquence}
{
Un fichier séquentiel impose un accès séquentiel à ses éléments.
D'autres structures de données peuvent être accédées
séquentiellement : par exemple le tableau et la liste mais aussi la
liste chainée (cf. votre cours de logique de 2ème année), la
communication sur un port série ou via le réseau, ... Les algorithmes
développés pour les fichiers peuvent être réutilisés pour
n'importe quelle séquence.}

{\sffamily\bfseries\upshape
Le tableau vu comme une séquence}

{
Un tableau permet un accès direct à ses éléments mais on peut tout aussi
bien y accéder en séquence. Comment traduire les notations introduites
pour les fichiers ? Il suffit d'utiliser un indice
courant et de le faire évoluer :}

{\sffamily
// Supposons tab : tableau [1 à n] de T et utilisons un indiceCourant}


\bigskip

{\sffamily
// ouvrir devient}

{\sffamily
\ \ indiceCourant {\textrm{[F0DF?]}} 0}


\bigskip

{\sffamily
// lire devient}

{\sffamily
\ \ indiceCourant {\textrm{[F0DF?]}}
indiceCourant + 1}

{\sffamily
\ \ \textstyleMotCl{si} indiceCourant ${\leq}$ N \textstyleMotCl{alors}}

{\sffamily
\ \ \ \ sauve {\textrm{[F0DF?]}} tab[
indiceCourant ]}

{\sffamily
\ \ \textstyleMotCl{fin} \textstyleMotCl{si}}


\bigskip

{\sffamily
//\ \ EOF devient}

{\sffamily
\ \ indiceCourant {\textgreater} N}

{\sffamily\bfseries\upshape
La liste vue comme une séquence}

{
Une liste aussi peut aussi être vue comme une séquence. Nous vous
laissons en exercice le soin de déterminer ce que deviennent les
primitives ouvrir/lire/\textsf{EOF} dans ce cas}

\section{Un accès différent aux fichiers/séquences}
{
Dans les fichiers comme nous les avons vus actuellement, une «~lecture~»
infructueuse est nécessaire pour se rendre compte
qu'il n'y a plus
d'informations à lire. }

{
On peut envisager des fichiers (ou plus généralement des séquences) où
on doit d'abord demander s'il reste
un élément à traiter avant de l'obtenir effectivement.
Ce schéma est celui qu'on trouve notamment dans }

\liststyleListv
\begin{itemize}
\item {
les curseurs en SQL}
\item {
les itérateurs en Java}
\item {
la lecture via la classe Scanner en Java}
\item {
la lecture dans certains langages de programmation comme Visual Basic}
\end{itemize}
{
On pourrait définir une telle séquence ainsi}

{\sffamily
\textstyleMotCl{classe} Séquence de T}

{\sffamily
\ \ \textstyleMotCl{méthode} débuter()\ \ \ \ }

{\sffamily
\ \ \textstyleMotCl{méthode} aEncore() {→}
booléen}

{\sffamily
\ \ \textstyleMotCl{méthode} suivant() {→} T}

{\sffamily
\textstyleMotCl{fin} \textstyleMotCl{classe}}

{
Pour illustrer ceci, écrivons l'algorithme qui parcourt
la séquence.}

{\sffamily
\textstyleMotCl{module} parcourir ( données : Séquence de T )}

{\sffamily
\ \ données.débuter()}

{\sffamily
\ \ \textstyleMotCl{tant} \textstyleMotCl{que} données.aEncore()
\textstyleMotCl{faire}}

{\sffamily
\ \ \ \ élément {\textrm{[F0DF?]}}
données.suivant()}

{\sffamily\itshape
\ \ \ \ // Traiter l'élément}

{\sffamily
\ \ \textstyleMotCl{fin} \textstyleMotCl{tant} \textstyleMotCl{que}\ \ }

{\sffamily
\textstyleMotCl{fin} \textstyleMotCl{module}}


\bigskip

{\sffamily\bfseries\scshape
Exercice – Adaptation aux tableaux et aux listes}

{
Que deviennent les méthodes de la classe Séquence de T pour un
tableau/une liste ?}

\section{}
\section{Exercices}
{
Pour ces exercices, vous pouvez utiliser les notations orienté objet ou
non mais restez cohérents au sein d'un même exercice
(pas de mélange).}

\liststyleExercice
\begin{enumerate}
\item {\sffamily\bfseries
Algorithmes de base}
\end{enumerate}
{
Soit le fichier FILE dont chacun des enregistrements est de type
\textbf{Elément}. Cette structure contient les champs
\textstyleCodeInsr{numéro} (entier), \textstyleCodeInsr{clé} (chaine)
et \ \textstyleCodeInsr{info} (chaine). Écrire des modules réalisant
les requêtes suivantes :}

\liststyleNumberingv
\begin{enumerate}
\item {
détermination de la \textbf{longueur} du fichier (c’est-à-dire le nombre
d’enregistrements du fichier)}
\item {
affichage de la valeur minimale du champ \textstyleCodeInsr{numéro}
parmi les enregistrements du fichier}
\item {
recherche du \textit{n}{}-ème enregistrement du fichier. L’entier
\textit{n} sera entré en paramètre et un paramètre booléen traduira
l’existence ou non de l’enregistrement recherché ; dans l’affirmative,
un troisième paramètre renverra cet enregistrement.}
\item {
recherche d’un élément connaissant sa «~clé~», c’est-à-dire renvoi en
paramètre du premier enregistrement du fichier dont le champ
\textstyleCodeInsr{clé} est égal à la valeur \textstyleCodeInsr{val}
entrée en paramètre. On renverra également un paramètre booléen
traduisant l’existence ou non de l’élément cherché. }
\end{enumerate}

\bigskip

\liststyleExercice
\begin{enumerate}
\item {\sffamily\bfseries
Copies et modifications de fichiers }
\end{enumerate}
{
Les énoncés suivants sont volontairement «~flous~». À vous de préciser
leur modalités de fonctionnement et de prévoir les différents scénarios
possibles.}

\liststyleNumberingv
\begin{enumerate}
\item {
Écrire un algorithme qui supprime le \textit{n}{}-ème élément d’un
fichier}
\item {
Écrire un algorithme qui supprime d’un fichier les éléments dont la
valeur d’un certain champ est égale à une valeur entrée en paramètre }
\item {
Écrire un algorithme qui insère dans un fichier une valeur entrée en
paramètre en \textit{n}{}-ème position}
\item {
Écrire un algorithme qui «~éclate~» un fichier en deux fichiers ; le
premier contiendra les éléments dont la valeur d’un certain champ est
inférieure à une valeur entrée en paramètre, et le deuxième contiendra
les autres éléments.}
\end{enumerate}

\bigskip

\liststyleExercice
\begin{enumerate}
\item {\sffamily\bfseries
Fichiers ordonnés}
\end{enumerate}
{
Soit FILENUM un fichier d’entiers. Écrire un algorithme :}

\liststyleNumberingv
\begin{enumerate}
\item {
qui vérifie si les éléments du fichier sont ordonnés par ordre
croissant}
\item {
qui recherche le mode d’ordonnancement du fichier (croissant,
strictement croissant, décroissant, strictement décroissant ou aucune
de ces possibilités)}
\end{enumerate}
{
En supposant à présent FILENUM ordonné en mode croissant, écrire un
algorithme :}

\liststyleNumberingv
\setcounter{saveenum}{\value{enumi}}
\begin{enumerate}
\setcounter{enumi}{\value{saveenum}}
\item {
qui insère la valeur VAL dans une copie de ce fichier}
\item {
qui crée une copie du fichier débarrassé de ses redondances}
\end{enumerate}

\bigskip

\liststyleExercice
\begin{enumerate}
\item {\sffamily\bfseries
Statistiques sur les profs}
\end{enumerate}
{
{Soit le fichier FILEPROFS des professeurs de la
HEB. La structure des enregistrements de ce fichier est
}{\textbf{Prof}}{
contenant les champs suivants :}}

{
nom : \ \ \ \ chaine}

{
âge : \ \ \ \ entier}

{
{sexe : \ \ \ \ logique
(}{\textbf{vrai}}{ si
masculin,
}{\textbf{faux}}{ si
féminin)}}

{
{dateIn : \ \ date d’entrée en fonction (objet
}\textstyleCodeInsr{{Date}}{)}}

{
Écrire l’algorithme qui :}

\liststyleNumberingv
\begin{enumerate}
\item {
affiche les noms des professeurs qui auront 20 ans d’ancienneté ou plus
au 31 décembre prochain}
\item {
affiche l’âge moyen du corps professoral}
\item {
affiche en \% la proportion de femmes dans le corps professoral}
\end{enumerate}

\bigskip

\liststyleExercice
\begin{enumerate}
\item {\sffamily\bfseries
La sélection des mannequins}
\end{enumerate}
{
Une agence de mannequins a organisé un concours ouvert aux débutantes
non professionnelles. Les postulantes inscrites ont toutes participé à
un examen de culture générale au cours duquel les questions étaient
rédigées indifféremment en français ou en néerlandais. Les
organisateurs du concours ont encodé tous les renseignements
nécessaires sur les candidates dans le fichier CANDI dont la structure
des enregistrements est de type \textbf{Candidat} comprenant les champs
suivants :}

{
nom : \ \ \ \ chaine}

{
adresse : \ \ chaine}

{
téléphone :\ \ chaine}

{
date : \ \ \ \ date de naissance
({objet}{\textbf{\textit{
}}}\textstyleCodeInsr{{Date}})}

{
taille : \ \ \ \ taille en mètre (réel)}

{
poids :\ \ \ \ poids en kilo (réel)}

{
pts : \ \ \ \ points obtenus à l’examen coté sur 150 (entier)}

{
En vue de réaliser une sélection, le jury a décidé d’éliminer les
candidates :}

\liststyleListv
\begin{itemize}
\item {
n’ayant pas atteint 18 ans au 1\textsuperscript{er} novembre de l’année
en cours ;}
\item {
ayant 30 ans ou plus au 1\textsuperscript{er} décembre de l’année en
cours ;}
\item {
de moins de 1m70 ou de plus de 1m90 ;}
\item {
ayant un IMC inférieur à 18 ou supérieur à 20, l’IMC (\textit{Indice de
Masse Corporelle}) étant le rapport entre le poids et le carré de la
taille ;}
\item {
ayant obtenu moins de 70\% à l’examen de culture générale.}
\end{itemize}
{
Écrire l’algorithme qui réalise, en une seule lecture du fichier CANDI,
cette sélection et qui écrit, sur le fichier SELECTION, les
enregistrements des candidates sélectionnées qui pourront continuer le
concours.}

\liststyleExercice
\begin{enumerate}
\item {\sffamily\bfseries
Le top 10}
\end{enumerate}
{
{Soit le fichier FILENB dont les enregistrements
sont des nombres entiers. Écrire} l’algorithme qui relève les dix plus
petites valeurs différentes dans ce fichier et les affiche.}

\liststyleExercice
\setcounter{saveenum}{\value{enumi}}
\begin{enumerate}
\setcounter{enumi}{\value{saveenum}}
\item[] 
\bigskip
\item {\sffamily\bfseries
{Le méli-mélo de
cartes}{. }}
\end{enumerate}
{
{Soit le fichier CARTES dont chaque
enregistrement est une structure
}{\textbf{Carte}}{
contenant les champs :}}

{
couleur : entier (0 pour PIQUE, 1 pour TREFLE, 2 pour CARREAU et 3 pour
CŒUR)~}

{
figure : \ \ entier (1 pour l’AS, 2 pour le DEUX,…, 10 pour le DIX, 11
pour le VALET,}

{
\ \ \ \ \ \ 12 pour la DAME et 13 pour le ROI)}

{
{Le fichier décrit un nombre indéterminé de
cartes provenant de paquets différents ; comme le nombre de cartes est
élevé, chaque carte apparait certainement plusieurs fois dans le
fichier. Écrire} l’algorithme qui dénombre et affiche le nombre de jeux
complets de 52 cartes contenu dans ce fichier.}

\liststyleExercice
\setcounter{saveenum}{\value{enumi}}
\begin{enumerate}
\setcounter{enumi}{\value{saveenum}}
\item {\sffamily\bfseries
La partie d’échec}
\end{enumerate}
{
Lorsqu’une partie d’échec est interrompue, on enregistre la position de
toutes les pièces sur l’échiquier de manière à pouvoir les replacer le
jour où la partie va reprendre. Le fichier de sauvegarde est JEU et
contient des enregistrements de type \textbf{Case}, chacun de ceux-ci
contenant les 3 champs :}

{
\ \ nl : \ \ numéro de ligne où se trouve la pièce (entier entre 1 et
8)}

{
\ \ nc : \ \ numéro de colonne où se trouve la pièce (entier entre 1 et
8)}

{
\ \ val : \ \ valeur de la pièce de type \textbf{Pièce}}


\bigskip

{
\textbf{Pièce} contient 2 champs: }

{
type:\ \  \ les valeurs sont ROI, REINE, CAVALIER, FOU, TOUR, PION}

{
couleur: les valeurs sont BLANC et NOIR}


\bigskip

{
Écrire l’algorithme parcourant le fichier JEU et permettant :}

\liststyleListv
\begin{itemize}
\item {
de replacer les pièces dans le tableau \textstyleCodeInsr{échiquier}
(tableau à deux dimensions 8 x 8 de \textbf{Pièce})}
\item {
de compter combien de cavaliers sont encore présents sur l’échiquier}
\item {
d’afficher la couleur du joueur possédant encore le plus de pièces sur
l’échiquier.}
\end{itemize}
\liststyleExercice
\begin{enumerate}
\item {\sffamily\bfseries
Les stages}
\end{enumerate}
{
Soit le tableau \textstyleCodeInsr{stages} de \textit{n} éléments de
structure \textbf{Stage} contenant les six champs : \ \ }

{
réfStage  : \ \ référence du stage}

{
âgeMin : \ \ âge minimum requis pour pouvoir participer}

{
âgeMax : \ \ âge maximum permis pour pouvoir participer}

{
nbMin : \ \ nombre minimum de participants requis pour que ce stage }

{
\ \ \ \ puisse avoir lieu}

{
nbMax : \ \ nombre maximum de participants admis}

{
nbDemande : \ \ nombre de demandes valables (au point de vue des limites
}

{
\ \ \ \ d’âges) déjà enregistrées avant aujourd’hui.}


\bigskip

{
Soit le fichier \textstyleCodeInsr{DEMANDES} des nouvelles demandes de
participation à un stage. La structure \textbf{Demande} d’un
enregistrement de ce fichier est la suivante :}


\bigskip


\bigskip

{
nom : \ \ \ \ nom et prénom du candidat}

{
âge : \ \ \ \ âge du candidat}

{
adresse : \ \ adresse du candidat}

{
réfStage : \ \ référence du stage pour lequel la demande est formulée.}


\bigskip

{
Écrire un module qui reçoit en paramètres le tableau
\textstyleCodeInsr{stages}, le tableau \textstyleCodeInsr{messages} de
6 éléments de type chaine (décrit ci-dessous) et le fichier
\textstyleCodeInsr{DEMANDES}. Ce module va traiter les nouvelles
demandes en mettant à jour le tableau \textstyleCodeInsr{stages} et en
créant un fichier \textstyleCodeInsr{REPONSES} de structure
d’enregistrement \textbf{InfoOut}. Cette structure comprend les champs
suivants :}


\bigskip

{
requête : \ \ structure \textbf{Demande}\textbf{\textit{ }}reprenant
toutes les informations }

{
\ \ \ \ fournies par un candidat}

{
msg : \ \ \ \ chaine de caractères devant contenir le message à envoyer
au}

{
\ \ \ \ demandeur, issu du tableau \textstyleCodeInsr{messages}}

{
Les messages possibles sont un des six suivants : }

\begin{center}
\tablehead{}
\begin{supertabular}{|m{0.36400002cm}|m{13.407001cm}|}
\hline
 1 &
 Le stage que vous avez demandé \textbf{n’existe
pas} ! La référence est incorrecte\\\hline
 2 &
 Votre êtes \textbf{trop jeune} pour participer
au stage que vous avez demandé !\\\hline
 3 &
 Vous être \textbf{trop âgé} pour participer au
stage que vous avez demandé !\\\hline
 4 &
 Désolé de ne pas prendre votre demande en
considération : le stage demandé \textbf{est complet}\\\hline
 5 &
 Votre inscription pour le stage est \textbf{mise
en attente} car il n’y a pas encore assez d’inscrits\\\hline
 6 &
 Votre inscription pour le stage demandé
\textbf{est acceptée}\\\hline
\end{supertabular}
\end{center}
{
\textstyleMotCl{\textmd{\textit{Notez qu’un seul message doit être
fourni en réponse à une demande. L’ordre de priorité des messages
correspond à l’ordre de leur apparition dans le tableau.}}}}

\liststyleExercice
\setcounter{saveenum}{\value{enumi}}
\begin{enumerate}
\setcounter{enumi}{\value{saveenum}}
\item {\sffamily\bfseries
Séquence. }
\end{enumerate}
{
Écrivez un module qui reçoit une séquence d'entiers et
qui}

\liststyleNumberingv
\begin{enumerate}
\item {
en affiche tous les éléments}
\item {
détermine la valeur la plus grande (erreur si la séquence est vide)}
\item {
détermine si la séquence est (strictement) croissante (erreur si la
séquence est vide)}
\item {
détermine si la valeur \textstyleCodeInsr{val} (passée en paramètre) est
présente ou non.}
\end{enumerate}
